% TOP_PROBLEM: {"id": 2834, "title": "Kirby Problem 3.36", "category_id": 7, "category": {"id": 7, "name": "topology", "display_name": "Topology", "description": "Properties preserved under continuous deformations.", "slug": "topology", "order_index": 7, "created_at": "2026-07-31T15:26:25.671Z"}, "status": "open", "problem_number": "KP-3.36", "set_id": 11}
% FIRST_POSTED: 2026-10-01
% ENTRY_KIND: statement_only
% UnsolvedMath ID 2834; source catalogue code KP-3.36
% !TeX program = lualatex
% Definition and sources only; this file is not an LLM solution attempt.
\documentclass[11pt,a4paper]{article}
\usepackage[margin=25mm]{geometry}
\usepackage{fontspec}
\setmainfont{lmroman10-regular.otf}[BoldFont=lmroman10-bold.otf,
 ItalicFont=lmroman10-italic.otf,BoldItalicFont=lmroman10-bolditalic.otf]
\usepackage{amsmath,amssymb,mathtools}
\usepackage{unicode-math}
\setmathfont{latinmodern-math.otf}
\AtBeginDocument{\let\setminus\smallsetminus}
\usepackage{xurl,hyperref}
\hypersetup{colorlinks=true,urlcolor=blue}
\setcounter{secnumdepth}{0}
\setlength{\parindent}{0pt}
\setlength{\parskip}{5pt}
\setlength{\emergencystretch}{3em}
\begin{document}
\section*{Kirby Problem 3.36}\label{2834}
{\small UnsolvedMath ID 2834 · KP-3.36 · Topology · Kirby's Problems in Low-Dimensional Topology}
\subsection{Definitions and mathematical statement}
Let $F$ be a connected surface and $M$ a three-manifold, and let
$f:F\to M$ be a smooth immersion. Its normal line bundle is
$\nu_f=f^*TM/TF$. The immersion is two-sided if this line bundle is trivial,
equivalently if it admits a continuous nowhere-zero normal vector field.
Choose basepoints to define $f_*:\pi_1(F)\to\pi_1(M)$; injectivity and the
existence of the loops below do not depend on that choice.

Assume that $f_*$ is not injective. The simple loop conjecture asks whether
there exists an embedded circle $c:S^1\hookrightarrow F$ with
\[
 [c]\ne1\text{ in }\pi_1(F),\qquad
 f_*[c]=1\text{ in }\pi_1(M).
\]
An unbased curve determines a conjugacy class, and lying in the kernel is
well-defined because the kernel is normal. The first condition says the
curve is essential, or not null-homotopic in $F$; without it, tiny circles
bounding disks would make the question vacuous.

The image loop $f\circ c$ need not be embedded in $M$. It is the curve on
the domain surface that must be simple. Thus the question asks whether a
nontrivial kernel necessarily contains a simple representative, not merely
whether some possibly self-intersecting loop lies in it.
\subsection{Short English statement}
If a two-sided immersed surface fails to inject on fundamental groups, must an essential simple loop on that surface become null-homotopic in the ambient three-manifold?
\subsection{Sources}
\begin{itemize}
\item[S1] ULAM.AI and the UnsolvedMath Contributors, supplied problems.json, record 2834 (KP-3.36), Kirby's Problems in Low-Dimensional Topology. Catalogue identity and source transcription; CC BY 4.0.
\url{https://huggingface.co/datasets/ulamai/UnsolvedMath}.
\item[S2] K3: A New Problem List in Low-Dimensional Topology, Problem 3.36. Expanded formulation of the supplied catalogue statement.
\url{https://math.berkeley.edu/sites/default/files/surv-295-ruberman-watermarked-author-pdf.pdf}.
\end{itemize}
\end{document}
